%%%%%%%%%%%%%%%%%%%%%%%%%%%%%%%%%%%%%%%%%%%%%%%%%%%%%%%%%%%%%%%%%%%%%%
% Lớp 8 | Chương IV. Định lí Thalès | Đường trung bình của tam giác
% Dạng 8H4?16-1: Tính độ dài nhờ đường trung bình
%%%%%%%%%%%%%%%%%%%%%%%%%%%%%%%%%%%%%%%%%%%%%%%%%%%%%%%%%%%%%%%%%%%%%%


% nguồn: dan_tu_AI (nhập 16/11/2026)
\begin{bt}%[8H4H16-1]%[ma:4b98e5ff]
	Cho tam giác $ABC$ có $BC=10\text{ cm}$. Gọi $M$, $N$ lần lượt là trung điểm của $AB$, $AC$. Tính độ dài $MN$.
	\loigiai{
		$M$, $N$ lần lượt là trung điểm của $AB$, $AC$ nên $MN$ là đường trung bình của tam giác $ABC$.\\
		Do đó $MN=\dfrac12BC=\dfrac12\cdot 10=5\text{ (cm)}$.
	}
\end{bt}

% nguồn: dan_tu_AI (nhập 16/11/2026)
\begin{bt}%[8H4H16-1]%[ma:91ff4a5c]
	Cho tam giác $ABC$ có $AB=8\text{ cm}$, $AC=10\text{ cm}$, $BC=12\text{ cm}$. Gọi $M$, $N$, $P$ lần lượt là trung điểm của $AB$, $AC$, $BC$. Tính chu vi tam giác $MNP$.
	\loigiai{
		$MN$, $NP$, $MP$ là các đường trung bình của tam giác $ABC$ nên
		\[ MN=\dfrac{BC}{2}=6,\quad NP=\dfrac{AB}{2}=4,\quad MP=\dfrac{AC}{2}=5\ \text{(cm)}. \]
		Chu vi tam giác $MNP$: $6+4+5=15\text{ (cm)}$.
	}
\end{bt}

% nguồn: dan_tu_AI (nhập 16/11/2026)
\begin{bt}%[8H4H16-1]%[ma:dafc2d87]
	Cho tam giác $ABC$ có $AC=12\text{ cm}$, $BC=9\text{ cm}$. Gọi $M$ là trung điểm của $AB$. Qua $M$ kẻ đường thẳng song song với $BC$ cắt $AC$ tại $N$. Tính độ dài $AN$ và $MN$.
	\loigiai{
		Tam giác $ABC$ có $M$ là trung điểm $AB$ và $MN\parallel BC$ nên $N$ là trung điểm $AC$.\\
		Do đó $AN=\dfrac{AC}{2}=6\text{ (cm)}$.\\
		$MN$ là đường trung bình của tam giác $ABC$ nên $MN=\dfrac{BC}{2}=4{,}5\text{ (cm)}$.
	}
\end{bt}

% nguồn: dan_tu_AI (nhập 16/11/2026)
\begin{bt}%[8H4H16-1]%[ma:27177ac0]
	Để đo khoảng cách giữa hai điểm $A$ và $B$ bị ngăn cách bởi một cái ao, người ta chọn điểm $C$ ở bên ngoài ao, rồi xác định trung điểm $M$ của $CA$ và trung điểm $N$ của $CB$. Đo được $MN=32\text{ m}$. Tính khoảng cách $AB$.
	\loigiai{
		$M$, $N$ lần lượt là trung điểm của $CA$, $CB$ nên $MN$ là đường trung bình của tam giác $CAB$.\\
		Do đó $MN=\dfrac12AB$, suy ra $AB=2\cdot MN=2\cdot 32=64\text{ (m)}$.
	}
\end{bt}

% nguồn: dan_tu_AI (nhập 16/11/2026)
\begin{bt}%[8H4H16-1]%[ma:97b2a24e]
	Cho tứ giác $ABCD$ có $AC=10\text{ cm}$, $BD=8\text{ cm}$. Gọi $E$, $F$, $G$, $H$ lần lượt là trung điểm của $AB$, $BC$, $CD$, $DA$. Tính chu vi tứ giác $EFGH$.
	\loigiai{
		$EF$ là đường trung bình của tam giác $ABC$ nên $EF=\dfrac{AC}{2}=5\text{ cm}$.\\
		$GH$ là đường trung bình của tam giác $ACD$ nên $GH=\dfrac{AC}{2}=5\text{ cm}$.\\
		$FG$ là đường trung bình của tam giác $BCD$ nên $FG=\dfrac{BD}{2}=4\text{ cm}$.\\
		$EH$ là đường trung bình của tam giác $ABD$ nên $EH=\dfrac{BD}{2}=4\text{ cm}$.\\
		Chu vi tứ giác $EFGH$: $5+4+5+4=18\text{ (cm)}$.
	}
\end{bt}

% nguồn: dan_tu_AI (nhập 16/11/2026)
\begin{bt}%[8H4V16-1]%[ma:077f2409]
	Cho tam giác $ABC$ có $AC=12\text{ cm}$, đường trung tuyến $AM$. Gọi $D$ là trung điểm của $AM$. Đường thẳng $BD$ cắt $AC$ tại $E$. Tính độ dài $AE$.
	\loigiai{
		Gọi $F$ là trung điểm của $EC$.\\
		Tam giác $BEC$ có $M$, $F$ lần lượt là trung điểm của $BC$, $EC$ nên $MF$ là đường trung bình, suy ra $MF\parallel BE$, tức $MF\parallel DE$.\\
		Tam giác $AMF$ có $D$ là trung điểm của $AM$ và $DE\parallel MF$ nên $E$ là trung điểm của $AF$, suy ra $AE=EF$.\\
		Mà $EF=FC$ nên $AE=EF=FC=\dfrac{AC}{3}=\dfrac{12}{3}=4\text{ (cm)}$.
	}
\end{bt}
